\uuid{zru3}
\titre{Implication/équivalence}
\auteur{bodin ; hanani ; mzari}
\organisation{exo7}
\licence{CC BY-NC-SA 4.0}
\datecreate{2020-01-20}
\niveau{L1}
\module{Algèbre}
\chapitre{Logique, ensemble, raisonnement}
\sousChapitre{Absurde et contraposée}
\difficulte{3}
\theme{Raisonnement, implication, équivalence}
\exo7qcmid{lille-1.logique.028}

\contenu{
\begin{question}
\qclassification{100.03,100.04}
\qtags{motcle=implication/équivalence}
Par quelle phrase puis-je remplacer la proposition logique "$P \Longleftarrow Q$" ?
\begin{answers}
    \good{"$P$ si $Q$"}

    \bad{"$P$ seulement si $Q$"}

    \bad{"$Q$ est une condition nécessaire pour obtenir $P$"}

    \good{"$Q$ est une condition suffisante pour obtenir $P$"}
\end{answers}
\begin{explanations}
C'est plus facile si on comprend que "$P \Longleftarrow Q$", c'est "$Q \implies P$", autrement dit "si $Q$ est vraie, alors $P$ est vraie".
\end{explanations}
\end{question}
}
